\section{Motivation}
\begin{frame}
\frametitle{Motivation}
\framesubtitle{Overall strategy}

\begin{enumerate}
  \item First thought: ``operator fusion is a well researched problem''
  \pause
  \begin{itemize}
    \item How do production-grade optimizers approach this? (XLA, TensorRT, TVM, etc.)
    \item They usually employ rule-based methods (e.g., epilogue-fuse ReLU after Matmul)
  \end{itemize}
  \pause
  \item Rule-based methods work well in practice
  \begin{itemize}
    \item An operator's performance characteristics are often associated with its kind
    \item e.g., ReLU is memory bound and (non-skinny) Matmul is compute bound
  \end{itemize}
  \pause
  \item However, in this challenge...
  \begin{itemize}
    \item There are only two kinds of operator (Matmul and Pointwise)
    \item Each operator can have arbitrary cost using \texttt{base\_cost}
    \item Fusion opportunities are unlimited\pause *
  \end{itemize}
\end{enumerate}
\pause
\begin{tcolorbox}
  Search-based method is needed for this challenge
\end{tcolorbox}
\end{frame}
